%!TEX program = pdflatex
% Mathematics article -- amsart, the neutral class of the subject. Drafting
% here and converting on acceptance costs nothing, because a journal class
% changes the front matter and the page layout and leaves the body alone.
\documentclass[11pt,reqno]{amsart}

% amsart's own page is set for the AMS printed page and reads narrow on
% screen; a journal class resets it at submission.
\usepackage[margin=1in]{geometry}

% amsart already loads amsmath, amsthm and amsfonts.
\usepackage{amssymb,mathtools}
\usepackage{graphicx}
\usepackage{bm}
\usepackage{booktabs}
\usepackage{xcolor}
\usepackage{hyperref}
% Links are coloured rather than boxed: the default frames print as
% rectangles around every citation and cross-reference.
\hypersetup{colorlinks=true,
            linkcolor=blue!55!black,
            citecolor=blue!55!black,
            urlcolor=blue!55!black}
\usepackage[capitalize,nameinlink]{cleveref}

% One shared counter, so a Lemma 2.3 and a Theorem 2.4 never collide.
\theoremstyle{plain}
\newtheorem{theorem}{Theorem}[section]
\newtheorem{lemma}[theorem]{Lemma}
\newtheorem{proposition}[theorem]{Proposition}
\newtheorem{corollary}[theorem]{Corollary}
\newtheorem{conjecture}[theorem]{Conjecture}

\theoremstyle{definition}
\newtheorem{definition}[theorem]{Definition}
\newtheorem{example}[theorem]{Example}

\theoremstyle{remark}
\newtheorem{remark}[theorem]{Remark}

% amsart numbers equations straight through; section numbering matches the
% theorem counter and survives inserting a section.
\numberwithin{equation}{section}

% A few lines -- the affiliation, the abstract, a theorem statement -- hold
% unbreakable runs slightly wider than the measure; let TeX stretch them
% rather than push type into the margin.
\emergencystretch=3em

\title{A Locally Finite Obstruction to Literal Fermionic Completeness}

\author{Zhiyuan Yao}
\address{Lanzhou Center for Theoretical Physics, Key Laboratory of Theoretical
  Physics of Gansu Province, Key Laboratory of Quantum Theory and Applications
  of MoE, Gansu Provincial Research Center for Basic Disciplines of Quantum
  Physics, Lanzhou University, Lanzhou 730000, China}
\email{yaozy@lzu.edu.cn}

\date{\today}

% AMS journals require a 2020 Mathematics Subject Classification, primary
% first; keywords are optional and printed with it.
% \subjclass[2020]{Primary 00A00; Secondary 00B00}
% \keywords{KEYWORDS}

\begin{document}

% amsart typesets the abstract from the title block, so it is
% declared inside the document and before that block.
\begin{abstract}
We distinguish the unrestricted inverse-limit space of charge-zero
restricted sine--Gordon towers from the algebraic graded direct sum.  Under
the literal componentwise definition, the former is strictly larger than the
finite module generated by fermionic coefficients and odd local integrals.
For odd power sums $p_{r,n}$, let
$H_{k,n}=\det(p_{2(i+j)+1,n})_{0\le i,j<2k+1}$.  Moment-matrix
factorization gives $H_{k,n}=0$ for $k\ge n$ and
$H_{k,k+1}\ne0$.  Thus
$K_n=M_{0,n}\sum_{k=0}^{n-1}H_{k,n}$ is componentwise polynomial and obeys
the annihilation recurrence, but contains nonzero shifted degrees
$8k^2+6k+1$ for all $k$.  Every finite fermionic module expression has
finite shifted-degree support, so $K$ is outside that module.  An independent
Hamel-dimension argument gives the same strict containment.  The obstruction
does not decide completeness after imposing finite-grade support.
\end{abstract}

\maketitle

\section{Introduction and category boundary}

The subject rests on the exactly solved lattice models and factorized
scattering theory of the nineteen-seventies and eighties
\cite{Baxter_Book,Zamolodchikov_1979fs}, and on the quantum inverse
scattering method that exhibited their common algebra
\cite{Sklyanin_1979qi,Kulish_1981yb}, out of which the integrable
structure of the sine--Gordon and related field theories grew.  Conformal field theory
supplied the second ingredient \cite{Belavin_1984ic}, with the structure of
its Virasoro modules --- the graded spaces a completeness statement
ultimately counts --- determined by Feigin and Fuchs \cite{Feigin_1984vm},
and with the massive integrable models obtained from it by perturbations that
retain infinitely many integrals of motion \cite{Zamolodchikov_1989if}.  The
integrable structure common to the two was developed as a quantum
Korteweg--de Vries theory in the series of Bazhanov, Lukyanov and
Zamolodchikov \cite{Bazhanov_1996is,Bazhanov_1997is,Bazhanov_1999is}, with
the form-factor programme for massive integrable models set out by Smirnov
\cite{Smirnov_Book} and the sine--Gordon expectation values and
thermodynamics developed alongside it
\cite{Lukyanov_1997ff,Lukyanov_1997ee,Destri_1995ua,Zamolodchikov_1990tb,Fateev_1999ev}.  Null-vector and
integrable-structure constructions supply the older field-theoretic setting
for the tower relations \cite{Babelon_1997nv}.

Fermionic creation operators then organized correlation functions and
descendant spaces across integrable lattice, conformal, and field theories,
in a programme running from the XXZ chain through its scaling limits and on
to conformal and sine--Gordon models
\cite{Boos_2005rq,Boos_2006ar,Boos_2007hg,Boos_2009,Jimbo_2010hg,JIMBO_2010oo,Boos_2010hg,Boos_2011fb,Negro_2013rr,Boos_2018nr}.  Jimbo, Miwa, and Smirnov formulated a
restricted sine--Gordon tower space and proposed generation by finitely many
fermionic coefficients over the polynomial ring of odd local integrals
\cite{Jimbo_2011fs}.  Subsequent work has obtained finite-level and special
parameter results \cite{Adler_2025fb,Adler_2026fb}.

The literal tower definition is componentwise: every particle-number
component is a Laurent polynomial satisfying its degree and recurrence
conditions.  It does not itself impose a global finite-grade bound.  This
distinction changes the completeness question from an algebraic direct sum
to an inverse product.  The construction below is locally finite in particle
number but has infinitely many global grades, which is what separates the
three spaces of \Cref{tab:spaces}.

\begin{table}[t]
\centering
\caption{Where the completeness question can be asked, and whether the tower
  $K$ constructed below lies there. Only the first row reads the definition
  literally.}
\label{tab:spaces}
\begin{tabular}{lll}
\toprule
Space & Global grade support & Contains the tower $K$ \\
\midrule
Literal componentwise tower space & unrestricted & yes \\
Algebraic graded direct sum & finite & no \\
Finite fermionic module & finite & no \\
\bottomrule
\end{tabular}
\end{table}

\begin{theorem}\label{thm:fermionic-obstruction}
Fix a generic irrational $\nu\in(0,1)$.  In the literal charge-zero
restricted tower space, define
\[
K_n=M_{0,n}\sum_{k=0}^{n-1}
\det\left(p_{2(i+j)+1,n}\right)_{0\le i,j<2k+1}.
\]
Then $K$ is an admitted nonzero tower but is not a finite combination, over
the polynomial ring generated by the odd local integrals, of fixed joint
Laurent coefficients of balanced fermionic creation-operator products.
This holds for either standard Laurent expansion chosen uniformly at zero or
at infinity.
\end{theorem}

\section{Tower objects and the missing global restriction}
\label{sec:towers}

Fix a generic irrational $\nu\in(0,1)$ throughout.

\begin{definition}
  \label{def:tower}
  For $l-n=c$, a \emph{charge-$c$ tower} is a family of Laurent-polynomial
  components
  \begin{equation*}
    L^{(l,n)}(S_1,\dots,S_l\mid B_1,\dots,B_{2n}),
  \end{equation*}
  antisymmetric in the $S$ variables and symmetric in the $B$ variables,
  subject to the componentwise degree condition
  $0\le\deg_{S_i}L^{(l,n)}\le2n-1$ and the annihilation recurrence
  \begin{equation}
    \label{eq:recurrence}
    \begin{aligned}
      &L^{(l,n)}(S_1,\dots,S_{l-1},B\mid B_1,\dots,B_{2n-2},B,-B)\\
      &\qquad=(-1)^cB\prod_{a=1}^{l-1}(B^2-S_a^2)\,
        L^{(l-1,n-1)}(S_1,\dots,S_{l-1}\mid B_1,\dots,B_{2n-2}).
    \end{aligned}
  \end{equation}
  Write $\mathcal T_{0,c}$ for the space of charge-$c$ towers.
\end{definition}

The primary charge-zero tower is
\begin{equation}
  \label{eq:primary}
  M_{0,n}(S_1,\dots,S_n)=\det\bigl(S_a^{2b-1}\bigr)_{1\le a,b\le n},
  \qquad M_{0,0}=1,
\end{equation}
in which the vacuum expectation value of the primary field, an overall
constant carried explicitly in \cite{Jimbo_2011fs} and evaluated in
\cite{Lukyanov_1997ee}, is set to
one; it multiplies every component alike and no statement below depends on it.
The odd local integrals act on components by multiplication by
\begin{equation}
  \label{eq:integrals}
  I_{2j-1,n}=\sum_{a=1}^{2n}B_a^{2j-1},
  \qquad
  \overline I_{2j-1,n}=\sum_{a=1}^{2n}B_a^{-(2j-1)} .
\end{equation}
These are the odd local integrals of motion of the model, whose construction
for quantum sine--Gordon and for perturbed conformal field theory goes back to
\cite{Sasaki_1988va,Eguchi_1989do,Zamolodchikov_1989if}; on a tower each acts
separately on every component, so the polynomial ring they generate acts by
multiplication and moves the joint degree by odd amounts.

With $P_n(S)=\prod_a(S-B_a)$, the creation operators are
\begin{equation}
  \label{eq:creation}
  \begin{aligned}
    (\psi_0^*(Z)L)^{(l+1,n)}
    &=\frac{\operatorname{Skew}_{S_0,\dots,S_l}
        \bigl(C_n(Z,S_0)L^{(l,n)}(S_1,\dots,S_l)\bigr)}{P_n(-Z)\,l!},\\
    (\chi_0^*(Z)L)^{(l-1,n)}
    &=\frac{L^{(l,n)}(Z,S_1,\dots)-L^{(l,n)}(-Z,S_1,\dots)}{2P_n(-Z)},
  \end{aligned}
\end{equation}
where $\operatorname{Skew}$ is unnormalized antisymmetrization in the
displayed $S$ variables,
\begin{equation*}
  \operatorname{Skew}_{S_0,\dots,S_l}f
  =\sum_{\sigma\in\mathfrak S_{l+1}}\operatorname{sgn}(\sigma)\,
    f(S_{\sigma(0)},\dots,S_{\sigma(l)}),
\end{equation*}
and the kernel is the signed expression
\begin{equation}
  \label{eq:kernel}
  C_n(S_1,S_2)=\frac{S_1}{4\nu}
    \sum_{\epsilon_1,\epsilon_2=\pm1}
    \frac{P_n(\epsilon_1S_1)P_n(\epsilon_2S_2)}
         {\epsilon_1S_1+\epsilon_2S_2},
\end{equation}
a polynomial of joint degree $4n$.

The statement at issue asserts that every $L\in\mathcal T_{0,0}$ is a
\emph{finite} combination, over the polynomial ring generated by
\eqref{eq:integrals}, of fixed joint formal Laurent coefficients of
balanced products of the operators \eqref{eq:creation} acting on $M_0$.

Two features of that formulation are what the counterexample uses, and
both are visible in the source. First, $\mathcal T_{0,c}$ is defined
componentwise in \cite{Jimbo_2011fs}: each particle-number
component is separately a Laurent polynomial of the stated degree obeying
\eqref{eq:recurrence}, and neither the definition nor the completeness
proposal imposes a uniform total $B$-degree or a finite-grade support across
$n$. Second, the denominator in \eqref{eq:creation} leaves the expansion
point under-specified; the argument below works for either standard Laurent
expansion, at zero or at infinity, chosen uniformly for the tower-valued
series. A mode exponent that depends on $n$ is not a coefficient of one
formal Laurent series, so that reading is not available.

\section{Odd-moment Hankel determinants}
\label{sec:hankel}

For positive odd $r$ write
\begin{equation}
  \label{eq:power-sums}
  p_{r,n}=\sum_{a=1}^{2n}B_a^r=I_{r,n},
\end{equation}
and for $k\ge0$ set $m=2k+1$ and
\begin{equation}
  \label{eq:hankel}
  H_{k,n}=\det\bigl(p_{2(i+j)+1,n}\bigr)_{0\le i,j<m}.
\end{equation}

\begin{lemma}
  \label{lem:stability}
  Every $H_{k,n}$ is stable under insertion of the annihilation pair:
  \begin{equation}
    \label{eq:pair}
    p_{r,n}(B_1,\dots,B_{2n-2},B,-B)=p_{r,n-1}(B_1,\dots,B_{2n-2})
  \end{equation}
  for odd $r$, hence $H_{k,n}(B_1,\dots,B_{2n-2},B,-B)=H_{k,n-1}$.
\end{lemma}

\begin{proof}
  For odd $r$ the two inserted variables contribute $B^r+(-B)^r=0$.
  Substituting into every entry of \eqref{eq:hankel} gives the second
  assertion.
\end{proof}

\begin{lemma}
  \label{lem:vanishing}
  $H_{k,n}=0$ whenever $k\ge n$, and $H_{k,k+1}\ne0$.
\end{lemma}

\begin{proof}
  Let $V$ be the $m\times2n$ matrix $V_{i,a}=B_a^{2i}$ and
  $D=\operatorname{diag}(B_1,\dots,B_{2n})$. The matrix in
  \eqref{eq:hankel} is $VDV^{\mathsf T}$. If $k\ge n$ then $m=2k+1>2n$, so
  the rank of $VDV^{\mathsf T}$ is at most $2n<m$ and the determinant
  vanishes.

  The threshold is sharp. At $n=k+1$ there are $m+1$ particle variables; set
  one of them to zero and choose the remaining $m$ with nonzero distinct
  squares. Then
  \begin{equation}
    \label{eq:witness}
    H_{k,k+1}
    =\Bigl(\prod_{a=1}^{m}B_a\Bigr)
     \Bigl(\prod_{1\le a<b\le m}(B_b^2-B_a^2)\Bigr)^2\ne0,
  \end{equation}
  so the polynomial $H_{k,k+1}$ is not identically zero. Alternatively,
  choosing all $m+1$ variables distinct and positive makes
  $VDV^{\mathsf T}$ positive definite, which exhibits a nonzero value inside
  $(\mathbb C^*)^{2k+2}$.
\end{proof}

\begin{lemma}
  \label{lem:degree}
  $H_{k,n}$ is homogeneous of degree
  \begin{equation}
    \label{eq:dk}
    d_k=2m^2-m=8k^2+6k+1,
  \end{equation}
  and $(d_k)_{k\ge0}$ is strictly increasing.
\end{lemma}

\begin{proof}
  Every monomial of the determinant \eqref{eq:hankel} is a product of $m$
  entries $p_{2(i+\sigma(i))+1,n}$ over a permutation $\sigma$ of
  $\{0,\dots,m-1\}$, so its degree is
  $\sum_{i=0}^{m-1}\bigl(2(i+\sigma(i))+1\bigr)$. Since
  $\sum_i(i+\sigma(i))=m(m-1)$ independently of $\sigma$, every monomial
  carries the same degree $2m(m-1)+m$, which is \eqref{eq:dk}. Strict
  monotonicity is immediate from $d_{k+1}-d_k=16k+14>0$.
\end{proof}

\section{An admitted locally finite tower}
\label{sec:K}

\begin{definition}
  \label{def:K}
  Put
  \begin{equation}
    \label{eq:K}
    F_0=0,
    \qquad
    F_n=\sum_{k=0}^{n-1}H_{k,n},
    \qquad
    K_n=M_{0,n}F_n .
  \end{equation}
\end{definition}

By \Cref{lem:vanishing}, \eqref{eq:K} is also the componentwise value of the
formal sum $M_{0,n}\sum_{k\ge0}H_{k,n}$, every component of which is an
ordinary finite polynomial. This is the sense in which $K$ is locally finite
and globally unbounded, and it is the whole construction.

\begin{proposition}
  \label{prop:admitted}
  $K\in\mathcal T_{0,0}$, and $K\ne0$.
\end{proposition}

\begin{proof}
  The primary determinant factorizes as
  \begin{equation}
    \label{eq:M-factor}
    M_{0,n}(S_1,\dots,S_{n-1},B)
    =B\prod_{a=1}^{n-1}(B^2-S_a^2)\,M_{0,n-1}(S_1,\dots,S_{n-1}),
  \end{equation}
  is alternating in $S$, and has $\deg_{S_i}\le2n-1$. The multiplier $F_n$
  is symmetric in $B$ and free of $S$, so $K_n$ has all the symmetries and
  degrees required by \Cref{def:tower}.

  Under insertion of the annihilation pair, \Cref{lem:stability,lem:vanishing}
  give
  \begin{equation}
    \label{eq:F-recurrence}
    F_n(B_1,\dots,B_{2n-2},B,-B)
    =\sum_{k=0}^{n-1}H_{k,n-1}
    =\sum_{k=0}^{n-2}H_{k,n-1}
    =F_{n-1},
  \end{equation}
  the removed term being $H_{n-1,n-1}=0$. Combining
  \eqref{eq:M-factor} and \eqref{eq:F-recurrence} gives
  \eqref{eq:recurrence} with $c=0$ for every $n$. Finally
  $K_1(S_1\mid B_1,B_2)=S_1(B_1+B_2)\ne0$.
\end{proof}

\begin{corollary}
  \label{cor:support}
  For every $k$ the component $K_{k+1}$ carries a nonzero homogeneous term
  $M_{0,k+1}H_{k,k+1}$, whose degree relative to the primary is $d_k$. Hence
  $K$ carries infinitely many distinct nonzero shifted degrees.
\end{corollary}

\begin{proof}
  Immediate from \eqref{eq:witness} and \Cref{lem:degree}, the degrees $d_k$
  being distinct by strict monotonicity.
\end{proof}

\section{Finite fermionic expressions have finite support}
\label{sec:finite-support}

\begin{lemma}
  \label{lem:baseline}
  Define the degree baseline
  \begin{equation}
    \label{eq:baseline}
    b(l,n)=2nl-n^2 .
  \end{equation}
  Then the primary \eqref{eq:primary} has degree $b(n,n)=n^2$, and the
  operators \eqref{eq:creation} shift the joint degree by exactly the
  baseline increments
  \begin{equation}
    \label{eq:increments}
    b(l+1,n)-b(l,n)=2n,
    \qquad
    b(l-1,n)-b(l,n)=-2n .
  \end{equation}
\end{lemma}

\begin{proof}
  The kernel $C_n$ has joint degree $4n$ by \eqref{eq:kernel} and
  $P_n(-Z)$ has degree $2n$, so $\psi_0^*(Z)$ raises the joint degree by
  $2n$; the numerator of $\chi_0^*(Z)$ leaves the degree unchanged while the
  same denominator lowers it by $2n$. These agree with
  \eqref{eq:increments}.
\end{proof}

\begin{proposition}
  \label{prop:finite-support}
  Every finite combination, over the polynomial ring generated by
  \eqref{eq:integrals}, of fixed joint Laurent coefficients of balanced
  products of \eqref{eq:creation} acting on $M_0$ has finitely many shifted
  degrees.
\end{proposition}

\begin{proof}
  Let $G(Z;S,B)=\sum_aZ^aG_a(S,B)$ be jointly homogeneous. Comparing the two
  sides after simultaneous scaling of all variables shows that each fixed
  coefficient $G_a$ is homogeneous, of the joint degree minus $a$. The same
  comparison applies to any finite list of variables $Z_i,X_i$ and to the
  expansion at zero or at infinity alike. With \Cref{lem:baseline}, every
  fixed joint coefficient of a balanced fermionic product acting on $M_0$
  therefore carries a single degree relative to \eqref{eq:baseline},
  independent of the particle number.

  Each $I_{2j-1}$ shifts that degree by $2j-1$ and each
  $\overline I_{2j-1}$ by $-(2j-1)$, and a polynomial has finitely many
  monomials. So a finite combination occupies finitely many shifted degrees.
\end{proof}

\section{The contradiction, and an independent dimension check}
\label{sec:contradiction}

\begin{proof}[Proof of \Cref{thm:fermionic-obstruction}]
  \Cref{prop:admitted} places $K$ in $\mathcal T_{0,0}$ and shows it is
  nonzero. \Cref{cor:support} gives it infinitely many distinct nonzero
  shifted degrees, while \Cref{prop:finite-support} allows every element of
  the proposed module only finitely many. Hence $K$ is not such a
  combination, for either uniform choice of expansion point, and for every
  admissible generic irrational $\nu$.
\end{proof}

The conclusion does not depend on the particular tower $K$, nor on the
choice of expansion point, as the following count shows independently.

\begin{proposition}
  \label{prop:dimension}
  $\mathcal T_{0,0}$ has uncountable Hamel dimension, whereas the proposed
  module is countably spanned. The containment is therefore strict.
\end{proposition}

\begin{proof}
  For a sequence $a=(a_0,a_1,\dots)\in\mathbb C^{\mathbb N}$ put
  \begin{equation}
    \label{eq:Ka}
    K^a_n=M_{0,n}\sum_{k=0}^{n-1}a_kH_{k,n},
  \end{equation}
  which is an admitted tower by the proof of \Cref{prop:admitted}. If $k$ is
  the first index with $a_k\ne0$, component $k+1$ carries the nonzero
  degree-$d_k$ term $a_kM_{0,k+1}H_{k,k+1}$, so $a\mapsto K^a$ is injective.
  The sequences $(1,\lambda,\lambda^2,\dots)$ for distinct
  $\lambda\in\mathbb C$ are linearly independent by the finite Vandermonde
  determinant, and there are uncountably many of them.

  The proposed module, by contrast, is spanned by a countable set: the
  balancing index $r$ is an integer, the fixed coefficient multi-indices lie
  in $\mathbb Z^{2r}$, and the monomials in the countable family of integral
  generators form a countable set.
\end{proof}

\section{The repaired conjecture, which this paper leaves open}
\label{sec:repair}

The natural repair replaces the unrestricted componentwise tower space by
its algebraic graded direct sum: require finite shifted-degree support, or
formulate completeness one homogeneous grade at a time. Either condition
excludes \eqref{eq:K} and \eqref{eq:Ka}, and either matches the reading
suggested by finite-level descendant spaces.

This paper neither proves nor disproves that repaired statement. The
completeness task is posed in \cite{Jimbo_2011fs}; with the local
integrals of motion acting nontrivially, the identification with a Virasoro
basis has been carried out only to level four \cite{Boos_2018nr}, later work
on the underlying expectation values does not raise that level
\cite{Adler_2025fb}, and the full construction of \cite{Adler_2026fb} is
confined to the free-fermion point. The intended finite-grade conjecture therefore remains
open.

\section{Conclusion}

Literal inverse-limit completeness fails because componentwise finiteness
does not imply finite global grade.  Requiring finite shifted-degree support
removes the constructed obstruction and matches the usual algebraic
descendant interpretation.  Whether the fermionic generators are complete
in that repaired finite-grade space remains open.

\section*{Acknowledgments}

Z.\ Y.\ acknowledges support from the Strategic Priority Research Program of the
Chinese Academy of Sciences under Grant No.~XDB1680102. Z.\ Y.\ also
acknowledges support by National Natural Science Foundation of China (Grant
No.~12247101). We also acknowledge Gewu Intelligence Lab for providing the
collaborative research environment in which this project was developed.

Large language models were used substantively in this work: OpenAI GPT-5.6 Sol
in solving the problem, including the formulation of arguments and the
derivation of intermediate steps, and Anthropic Claude Opus~5 in surveying the
literature, in checking every attribution against the cited sources, and in
drafting and revising the manuscript. The author specified the problem, guided
the approach to be taken, checked every argument and result they produced, and
takes full responsibility for the content of this manuscript.

\bibliographystyle{amsplain}
\bibliography{references,unverified}
\end{document}
